% =====================================================================
%  Overleaf: Menu -> Compiler.  File nay chay duoc voi CA XeLaTeX LAN pdfLaTeX
%  (khoi iftex ben duoi tu chon goi font phu hop). Neu van loi font tieng Viet
%  thi chon XeLaTeX cho chac.
%
%  Thu muc:   bao_cao.tex
%             hinh/*.pdf     <- copy cac file PDF notebook sinh ra vao day
%
%  Thieu hinh nao thi macro \hinh ve mot o trong thay cho no -> van compile duoc.
% =====================================================================
\documentclass[aspectratio=169, 11pt]{beamer}

\usepackage{iftex}
\ifPDFTeX
  \usepackage[utf8]{inputenc}
  \usepackage[T5]{fontenc}      % T5 = bang ma tieng Viet (goi vntex, co san tren Overleaf)
  \usepackage{lmodern}
\else
  \usepackage{fontspec}
  \setsansfont{Latin Modern Sans}   % font nay LUON co trong TeX Live
\fi

\usepackage{amsmath, amssymb}
\usepackage{booktabs}
\usepackage{graphicx}
\usepackage{xcolor}

\usetheme{default}
\usecolortheme{seahorse}
\setbeamertemplate{navigation symbols}{}
\setbeamertemplate{footline}[frame number]
% mau giu nguyen nhu ban slide cu: structure = xanh navy tram
\setbeamercolor{structure}{fg=blue!55!black}
\colorlet{primary}{blue!55!black}     % \textcolor{primary} van dung duoc, cung mot mau
\setbeamerfont{frametitle}{size=\large, series=\bfseries}

% Tim hinh o CA HAI cho: hinh/ten.pdf va ten.pdf ngay canh file .tex.
% Tren Overleaf rat hay xay ra viec upload file phang ra goc thay vi vao thu muc,
% luc do \IfFileExists{hinh/...} se fail va chi ve ra o trong.
\graphicspath{{hinh/}{./}}
\newcommand{\vehinh}[2]{%
  \IfFileExists{hinh/#2}%
    {\includegraphics[width=\linewidth, height=#1\textheight, keepaspectratio]{hinh/#2}}%
    {\IfFileExists{#2}%
      {\includegraphics[width=\linewidth, height=#1\textheight, keepaspectratio]{#2}}%
      {\fbox{\parbox[c][#1\textheight][c]{0.93\linewidth}{\centering\footnotesize\ttfamily
        #2\\[2pt]\rmfamily\textcolor{red}{không tìm thấy \texttt{hinh/#2} lẫn \texttt{./#2}}}}}}}
% \hinh  -> cao toi da 0.34 khung (khi xep CHONG hai hinh trong mot cot)
% \hinhto -> cao toi da 0.60 khung (khi moi cot chi co MOT hinh)
\newcommand{\hinh}[1]{\vehinh{0.34}{#1}}
\newcommand{\hinhto}[1]{\vehinh{0.60}{#1}}

\title{Tối ưu hoá hàm log-loss bằng các thuật toán gradient}
\subtitle{Santander Product Recommendation}
\author{Nhóm \ldots}
\date{\today}

\begin{document}

\frame{\titlepage}

% =====================================================================
\begin{frame}{Giới thiệu Dataset \& Bài toán}

\textbf{Bài toán:} dự đoán bán chéo (cross-sell) sản phẩm tài chính $\Rightarrow$ phân loại nhị
phân bằng Logistic Regression, tối thiểu hoá log-loss \textbf{có regularization}:
\[
  F(\mathbf{w})=\underbrace{\frac1n\sum_{i=1}^{n}\Big[\log\big(1+e^{z_i}\big)-y_i z_i\Big]}_{\ell(\mathbf{w})\ \text{(log-loss)}}
  \;+\;\lambda R(\mathbf{w}),\qquad \mathbf{z}=X\mathbf{w}
\]
\[
  R(\mathbf{w})=\tfrac12\|\mathbf{w}\|_2^2\ \text{(Ridge/L2)}
  \qquad\qquad
  R(\mathbf{w})=\|\mathbf{w}\|_1\ \text{(Lasso/L1)}
\]
\textbf{Dataset:} Santander Product Recommendation (Kaggle) — $\sim 7{,}19$ triệu bản ghi khách hàng theo tháng.
Trong đó có 1.842.484 quan sát được dùng để huấn luyện.
\begin{center}
\small
\begin{tabular}{@{}lcl@{}}
\toprule
\textbf{Nhóm đặc trưng} & \textbf{Số} & \textbf{Mô tả} \\
\midrule
LAG & 24 & Sở hữu sản phẩm tháng trước (\texttt{ind\_*\_ult1\_lag}) \\
Standardized & 6  & \texttt{age}, \texttt{antiguedad}, \texttt{renta}, \texttt{ind\_nuevo}, \dots \\
Categorical & 7  & Giới tính, phân khúc, kênh, quan hệ (dummies) \\
\midrule
\textbf{Tổng} & \textbf{37} & Đầu vào mô hình \\
\bottomrule
\end{tabular}
\end{center}

\vspace{-0.2em}
{\footnotesize \textit{Target:} sở hữu sản phẩm ở tháng hiện tại (\texttt{ind\_*\_ult1}), suy ra từ nhóm SP (24 sản phẩm).}

\end{frame}

% =====================================================================
\begin{frame}[t]{Bộ tham số dùng để chạy}
\footnotesize
Bộ 1 là bộ đã chạy để lấy các hình phía sau. Bộ 2 là bộ mở rộng, dùng để kiểm chứng lại
kết luận: lưới rộng hơn nên chạm được cả hai đầu hỏng, và tăng $\lambda$ để thấy rõ tác dụng
của hiệu chỉnh.

\smallskip
\begin{center}
\begin{tabular}{@{}llll@{}}
\toprule
\textbf{Thuật toán} & \textbf{Tham số quét} & \textbf{Bộ 1 (đã chạy)} & \textbf{Bộ 2 (mở rộng)} \\
\midrule
GD bước cố định & $\alpha$ (bội của $1/L$) & $0.1;\,0.5;\,1;\,1.5;\,2$ & $0.5;\,1;\,2;\,5;\,20$ \\
GD backtracking & $\beta$ (hệ số co)       & $0.1;\,0.3;\,0.5;\,0.7;\,0.9$ & $0.2;\,0.5;\,0.8$ \\
                & $c$ (Armijo)             & $10^{-4}$ & $10^{-4};\,10^{-2}$ \\
SGD             & $\alpha_0$               & $0.01;\,0.05;\,0.2;\,1;\,5$ & $0.2;\,1;\,5;\,20;\,50$ \\
                & $B$ (batch)              & $4096$ & $256;\,4096$ \\
                & lịch teo bước            & cố định, $\alpha_0/\sqrt{k}$, $\alpha_0/k$ & như bộ 1 \\
Subgradient     & $\alpha_0$               & $0.5;\,2;\,8;\,30;\,100$ & $1;\,4;\,16;\,64$ \\
\midrule
\multicolumn{2}{@{}l}{Số vòng lặp / số epoch} & $250$ / $42$ & $100$ / $100$ \\
\multicolumn{2}{@{}l}{$\lambda$ (Ridge) / $\lambda_1$ (Lasso)} & $10^{-3}$ / $10^{-3}$ & $10^{-2}$ / $10^{-2}$ \\
\bottomrule
\end{tabular}
\end{center}

\smallskip
Nguyên tắc chọn lưới: phải \textbf{kẹp được nghiệm tối ưu ở giữa}. Nếu giá trị tốt nhất rơi
vào đúng hai đầu lưới thì chưa kết luận được, phải nới lưới ra rồi chạy lại. Bộ 1 có
$\alpha=2/L$ ở đầu trên và $\alpha_0=100$ ở subgradient chính là để chạm vào vùng hỏng.
\end{frame}
\begin{frame}[t]{GD với bước cố định: chọn $\alpha$ thế nào}
\begin{columns}[T]
\begin{column}{0.46\textwidth}\footnotesize
\[ \mathbf{w}_{k+1}=\mathbf{w}_k-\alpha\nabla F(\mathbf{w}_k) \]
Bước an toàn theo lý thuyết:
\[ \alpha=\frac1L,\qquad L=\frac{\sigma_{\max}(X)^2}{4n}+\lambda \]
(vì $\sigma'(z)\le\tfrac14$). Thử $\alpha\in\{0.1;0.5;1;1.5;2\}\times\tfrac1L$.

\smallskip
\textbf{Kinh nghiệm:}
\begin{itemize}\setlength\itemsep{0.2em}
  \item $\alpha=1/L$ chắc chắn hội tụ nhưng quá dè dặt. $L$ suy ra từ chặn
        $\sigma'(z)\le 1/4$ cho mọi điểm, trong khi tỉ lệ dương chỉ khoảng 1\% nên độ cong
        thực tế nhỏ hơn nhiều.
  \item Tăng $\alpha$ thì hội tụ nhanh hơn, nhưng chỉ đến ngưỡng $2/L$. Vượt qua đó thì
        $F$ dao động rồi phân kỳ.
  \item Phải dò bằng tay, và muốn có $1/L$ thì phải tính được $L$ trước.
\end{itemize}
\end{column}

\begin{column}{0.52\textwidth}
\hinh{gd_codinh_vonglap.pdf}\\[0.4em]
\hinh{gd_codinh_thoigian.pdf}
\end{column}
\end{columns}
\end{frame}

% =====================================================================
\begin{frame}[t]{GD với backtracking (Armijo)}
\begin{columns}[T]
\begin{column}{0.46\textwidth}\footnotesize
Mỗi vòng thử bước \textbf{gấp đôi} lần trước, chưa thoả thì nhân với $\beta$:
\[ F(\mathbf{w}-\alpha\nabla F)\le F(\mathbf{w})-c\,\alpha\,\|\nabla F\|^2 \]
với $c=10^{-4}$. Thử $\beta\in\{0.1;0.3;0.5;0.7;0.9\}$.

\smallskip
\textcolor{red}{\textbf{Chỗ dễ sai:}} mỗi vòng phải cho bước lớn dần lên. Nếu chỉ biết
co xuống từ $1/L$ thì backtracking không làm gì cả, vì $1/L$ vốn đã thoả điều kiện sẵn.

\smallskip
\textbf{Kinh nghiệm:}
\begin{itemize}\setlength\itemsep{0.2em}
  \item Xét theo số vòng lặp, năm giá trị $\beta$ cho đường gần như trùng nhau.
  \item Xét theo thời gian thì tách hẳn ra. $\beta$ lớn dò kỹ hơn nên mỗi vòng phải tính
        $F$ nhiều lần, thành ra chậm hơn.
  \item Một cấu hình có thể thắng ở hình này và thua ở hình kia, nên phải vẽ cả hai.
\end{itemize}
\end{column}

\begin{column}{0.52\textwidth}
\hinh{gd_bt_vonglap.pdf}\\[0.4em]
\hinh{gd_bt_thoigian.pdf}
\end{column}
\end{columns}
\end{frame}

% =====================================================================
\begin{frame}[t]{SGD: mini-batch và lịch teo bước}
\begin{columns}[T]
\begin{column}{0.46\textwidth}\footnotesize
\[ \nabla\ell(\mathbf{w})\approx\frac1B\sum_{i\in\mathcal{B}}\nabla\ell_i(\mathbf{w}) \]
Ước lượng \textbf{không chệch}, phương sai $\sigma^2/B$, rẻ hơn $n/B$ lần mỗi bước.
Điều kiện Robbins--Monro:
\[ \sum_k\alpha_k=\infty,\qquad \sum_k\alpha_k^2<\infty \]

\smallskip
\textbf{Kinh nghiệm:}
\begin{itemize}\setlength\itemsep{0.2em}
  \item Không dùng được backtracking cho SGD. Điều kiện Armijo cần so $F$ trước và sau,
        mà $F$ trên mini-batch là số nhiễu nên nó chấp nhận hay từ chối đều không đáng tin.
  \item Ở đây $d=38$ nên $B=1$ còn chậm hơn GD tính theo giây: một bước GD là một lệnh
        BLAS trên ma trận lớn, còn SGD là hàng nghìn lệnh rất nhỏ. Đo được một epoch tốn
        gấp 4 đến 8 lần một vòng GD, nên chọn $B=4096$.
  \item Teo bước chỉ có lợi khi chạy đủ lâu. Với 42 epoch, $\alpha_0/k$ teo quá nhanh nên
        chưa kịp xuống tới đáy.
\end{itemize}
\end{column}

\begin{column}{0.52\textwidth}
\hinh{sgd_eta_epoch.pdf}\\[0.4em]
\hinh{sgd_lich_thoigian.pdf}
\end{column}
\end{columns}
\end{frame}

% =====================================================================
\begin{frame}[t]{Subgradient cho Lasso: kết quả}
\begin{columns}[T]
\begin{column}{0.46\textwidth}\footnotesize
Quét $\alpha_0\in\{0.5;\,2;\,8;\,30;\,100\}$, lịch teo $\alpha_0/\sqrt{k}$, $\lambda_1=10^{-3}$,
250 vòng.

\smallskip
\begin{center}\scriptsize
\begin{tabular}{@{}rrc@{}}
\toprule
$\alpha_0$ & $F$ cuối & Hệ số $\neq0$ \\
\midrule
$0.5$ & 0.1\ldots & \ldots/37 \\
$2$   & 0.1\ldots & \ldots/37 \\
$8$   & \textbf{0.1\ldots} & \ldots/37 \\
$30$  & 0.1\ldots & \ldots/37 \\
$100$ & 0.1\ldots & \ldots/37 \\
\midrule
sklearn L1 & 0.1\ldots & \ldots/37 \\
\bottomrule
\end{tabular}
\end{center}

\smallskip
\textbf{Đọc được gì từ kết quả:}
\begin{itemize}\setlength\itemsep{0.2em}
  \item Đáy nằm ở $\alpha_0=8$, hai đầu lưới đều tệ hơn. Giá trị chọn được nằm giữa lưới
        chứ không phải ở biên, nên kết luận này dùng được.
  \item Độ thưa \textbf{không phụ thuộc $\alpha_0$}: cả năm lần chạy cho cùng số hệ số khác 0.
        Toạ độ $j$ chỉ đứng yên tại 0 khi $|\nabla_j\ell|\le\lambda_1$, mà điều kiện đó
        chỉ phụ thuộc dữ liệu và $\lambda_1$.
  \item $\alpha_0=100$ làm $F$ vọt lên trên rồi mới bò xuống suốt 200 vòng. Đây là bằng chứng
        trực tiếp cho việc subgradient không bảo đảm giảm ở mỗi bước.
  \item So với \texttt{liblinear}: $F$ chênh nhau ở chữ số thứ 6, nhưng code tự viết chậm
        hơn vài lần.
\end{itemize}
\end{column}

\begin{column}{0.52\textwidth}
\hinh{subgrad_vonglap.pdf}\\[0.4em]
\hinh{subgrad_thoigian.pdf}
\end{column}
\end{columns}
\end{frame}

% =====================================================================
\begin{frame}{So sánh: ba setup tốt nhất với sklearn}
\begin{columns}[T]
\begin{column}{0.5\textwidth}\hinhto{tongket_vonglap.pdf}\end{column}
\begin{column}{0.5\textwidth}\hinhto{tongket_thoigian.pdf}\end{column}
\end{columns}

\vspace{0.4em}
\begin{center}\small
\begin{tabular}{@{}lrr@{}}
\toprule
\textbf{Setup tốt nhất (Ridge)} & \textbf{$F$ cuối} & \textbf{Giây} \\
\midrule
GD bước cố định $\alpha=\ldots$                     & 0.0\ldots & \ldots \\
GD backtracking $\beta=\ldots$                 & 0.0\ldots & \ldots \\
SGD $\alpha_0=\ldots$, $B=4096$                  & 0.0\ldots & \ldots \\
\midrule
sklearn \texttt{LogisticRegression} (mặc định) & 0.0\ldots & \ldots \\
\bottomrule
\end{tabular}
\end{center}

\vspace{-0.1em}
{\footnotesize Điều kiện để so được: cả ba cùng tối thiểu hoá \emph{một} hàm $F$ (Ridge, cùng
$\lambda$). Quy đổi cho sklearn: nó dùng $\frac12\|\mathbf{w}\|^2+C\sum_i\ell_i$ (tổng), ta dùng
trung bình $\Rightarrow C=\dfrac{1}{n\lambda}$.}
\end{frame}

% =====================================================================
\begin{frame}{Lasso: subgradient so với thư viện}
\begin{columns}[T]
\begin{column}{0.5\textwidth}\hinhto{lasso_vonglap.pdf}\end{column}
\begin{column}{0.5\textwidth}\hinhto{lasso_thoigian.pdf}\end{column}
\end{columns}

\vspace{0.4em}
\begin{center}\small
\begin{tabular}{@{}lrrr@{}}
\toprule
\textbf{Lasso} & \textbf{$F$ cuối} & \textbf{Giây} & \textbf{Hệ số $\neq0$} \\
\midrule
Subgradient $\alpha_0=\ldots$                        & 0.1\ldots & \ldots & \ldots/37 \\
sklearn \texttt{liblinear} (\texttt{penalty='l1'}) & 0.1\ldots & \ldots & \ldots/37 \\
\bottomrule
\end{tabular}
\end{center}

\vspace{-0.1em}
{\footnotesize Để riêng vì subgradient tối thiểu hoá \emph{hàm khác}
($\lambda_1\|\mathbf{w}\|_1$ thay cho $\frac\lambda2\|\mathbf{w}\|_2^2$), nên đặt chung
một trục với Ridge là so sai.}
\end{frame}

% =====================================================================
\begin{frame}{Kết luận}
\begin{enumerate}\setlength\itemsep{0.45em}
  \item \textbf{Về giá trị hàm mục tiêu:} code tự viết đạt $F$ \textit{[thấp hơn / ngang]}
        sklearn ở chữ số thứ \ldots{}, tức phần cài đặt không có sai sót.
  \item \textbf{Về thời gian:} chậm hơn sklearn \ldots{} lần. sklearn dùng L-BFGS, một phương
        pháp tựa Newton khai thác thông tin độ cong nên cần rất ít vòng lặp.
  \item \textbf{Chọn tham số:} bước cố định phải dò tay; backtracking tự tìm bước mà không cần
        biết $L$, cái giá là vài lần tính $F$ thêm mỗi vòng.
  \item Số vòng lặp và thời gian là hai chuyện khác nhau. Một cấu hình có thể thắng ở
        hình này và thua ở hình kia, nên mọi kết luận đều phải đọc trên cả hai.
  \item \textbf{Điều kiện của bài toán:} với $\lambda=10^{-3}$ thì $L/\mu\approx1500$, nên GD cần
        $O(10^4)$ vòng để hội tụ chặt. Đây chính là lý do các phương pháp gia tốc
        ($O(\sqrt{L/\mu})$) và Newton (bậc hai) tồn tại.
\end{enumerate}

\vfill
\begin{center}\textcolor{primary}{\textbf{Cảm ơn --- Hỏi \& Đáp}}\end{center}
\end{frame}




\end{document}
